\section{A minimal type and its geometry}
\label{geom:section-geometry}

We continue with the fixed AEC $\K$ that is categorical in unboundedly
many cardinals.  Throughout this Section, a level is eligible for the
single-model working scheme and for the refinements of the calculus
that are explicitly invoked.  The first level chosen below is denoted
by $\lambda_0$.  The level used for all finite marked constructions will
be a later cardinal $\lambda$; these two choices serve different purposes.

\begin{lemma}[Preservation of nonalgebraicity]
\label{geom:nonalgebraic}
Let $p$ be a nonalgebraic singleton type over an object $A$ of a working
scheme.  Its free lift over an extension $B$ is nonalgebraic.  The same
assertion holds for a mixed arrow whenever the input arity and the
testing bounds lie below the smaller level.
\end{lemma}

\begin{proof}
Suppose that the free lift were realized by $c\in B$.  Include equality
with $c$ among its tests.  To evaluate that test, the free-extension
construction moves $c$ into $A$, over enough of the chosen
approximations for $p$.  Its value is then the value of equality with an
element of $A$.  Every such equality has value false for $p$, whereas
equality with its alleged realization $c$ has value true.  This
contradicts the evaluation rule in Theorem~\ref{calc:free-extension}.
The mixed-level construction uses the same rule with the bounds taken
below the input level.
\end{proof}

\begin{proposition}[A minimal type]
\label{geom:minimal-type}
There is a sufficiently closed single-model level $\lambda_0$, a model
$M_0\in\K$ of size $\lambda_0$, and a nonalgebraic singleton type
$p_0$ over $M_0$ such that, over every $M_0\leq_\K A\in\K$ of size
$\lambda_0$, there is exactly one nonalgebraic extension of $p_0$.
This extension is its free lift.
\end{proposition}

\begin{proof}
Choose $\lambda_0$ with the single-model working package, the
stationary calculus, the all-dominated hull construction, and the
singleton type bound of Theorem~\ref{lev:amalgamation}(3).
Thus there are at most $\lambda_0$ singleton
types over a model of that size.  We may require at the outset strict
closure under all the fixed bounds used below for finite tuples and
their bounded tests.

There is a nonalgebraic singleton type at this level.  For example,
extend an order presentation by a fresh skeleton point.  That point is
outside the old model: if it were a labeled term on finitely many old
indices, adding a second distinct skeleton point with the same order
pattern over those indices would make both equal that term, contrary
to the distinctness of the skeleton.

Suppose that no nonalgebraic singleton type at this level is minimal in
the stated sense.  A nonalgebraic type always has a nonalgebraic free
extension by Lemma~\ref{geom:nonalgebraic}.  Consequently every such
type has two distinct nonalgebraic extensions over some model of size
$\lambda_0$.

Let $\epsilon$ be the least infinite cardinal with
$2^\epsilon>\lambda_0$.  Then $\epsilon\leq\lambda_0$ and
$2^{|i|}\leq\lambda_0$ for $i<\epsilon$.  Construct a continuous chain
$(A_i:i\leq\epsilon)$ of models of size $\lambda_0$ and types
$p_\eta$ over $A_i$, for $\eta\in{}^i2$, as follows.  At a successor
step, split every current $p_\eta$ over a model extending $A_i$.
There are at most $\lambda_0$ such models.  Successive amalgamation
over $A_i$, followed by unions, puts all of them into one next base
$A_{i+1}$ of size $\lambda_0$.  Freely extend each of the two split
types to that base.  Their restrictions to the splitting model are
different, so they remain different, and
Lemma~\ref{geom:nonalgebraic} keeps them nonalgebraic.

Along each branch retain a witness model and the same distinguished
realization.  The required successor map exists by amalgamation with
identification of realizations having the same type over the preceding
base.  At a limit take the unions of the common bases and, separately
on each branch, of its coherent witness models.  The realization is
outside every base stage, so it is outside the union base.  All these
unions have size $\lambda_0$ because the construction has length at
most $\lambda_0$.  At the final stage distinct branches have different
types over $A_\epsilon$, witnessed at their first splitting stage.
This gives more than $\lambda_0$ singleton types over $A_\epsilon$,
contrary to the type bound.  A minimal $p_0$ therefore exists.  Its
unique nonalgebraic extensions are free by existence of free lifts.
\end{proof}

The passage from a minimal type to closure and dimension follows the
geometric pattern of first-order categoricity developed through strongly
minimal sets; see Baldwin and Lachlan \cite[Sections~1--2]{BaldwinLachlan1971}.
Here the required containment, exchange, and change-of-base properties
are established directly for the AEC and its locally constructed types.

\subsection{Finite configurations over the chosen model}

\begin{definition}
\label{geom:points}
For $M_0\leq_\K M$ with $\|M\|\geq\lambda_0$, let $P(M)$ be the
set of realizations of $p_0$ in $M$.  More explicitly, $a\in P(M)$
means that in a strong section $N\leq_\K M$ of size $\lambda_0$
containing $M_0\cup\{a\}$, the type $\tp(a/M_0)$ is $p_0$.
A finite ordered list of distinct points in $P(M)$ is independent if
its type over $M_0$, computed in such a section, is the product of the
corresponding copies of $p_0$.
\end{definition}

These definitions do not depend on the section.  Two sections lie in
a third section of size $\lambda_0$, by the L\"owenheim--Skolem axiom
and coherence.  The embeddings into that section preserve their types
over $M_0$.  Products are permutation invariant, so independence is a
property of the underlying finite set.  Strong embeddings over $M_0$
preserve and reflect both membership in $P$ and finite independence.

We give a precise meaning to the containment test used in constructing
closure.  Fix a finite configuration $F\cup\{x\}\subseteq P(M)$ and
a size-$\lambda_0$ section $N$ containing it and $M_0$.  Say that $F$
\emph{forces} $x$ if, whenever $f:N\to W$ is a strong embedding fixing
$M_0$, with $W$ of size $\lambda_0$, and
$M_0\leq_\K D\leq_\K W$ has size $\lambda_0$ and contains $f[F]$,
one has $f(x)\in D$.  Thus the test permits a common extension of the
finite configuration and the proposed model; it does not presume that
intersections of strong models are strong models.

The test is invariant under the type of the entire finite
configuration over $M_0$.  Indeed, an identification of two copies of
that type is witnessed in an amalgam.  Any further configuration
witnessing failure of forced containment can be amalgamated with that
identification over the first witness model.  Its avoiding submodel
and the omitted point remain disjoint at that point because the maps
are injective.  This transports the failure to the second copy.  The
reverse transport is identical with the copies interchanged.

\begin{lemma}[The containment test]
\label{geom:containment-test}
For a finite independent set $S\subseteq P(M)$ and $x\in P(M)$,
$S$ fails to force $x$ if and only if $S\cup\{x\}$ is independent
with $x\notin S$.
\end{lemma}

\begin{proof}
First suppose that a coupling contains a model $D$ over $M_0$ and $S$
but omits $x$.  The type of $x$ over $D$ is a nonalgebraic extension of
$p_0$.  By Proposition~\ref{geom:minimal-type} it is the free lift of
$p_0$.  Restricting its parameters to $S$ shows, by
Proposition~\ref{calc:products}, that $x$ has product type with $S$ over
$M_0$.

Conversely suppose that $S,x$ has product type.  Use
Theorem~\ref{iso:all-lengths} to obtain a model all-dominated over
$M_0+S$, and amalgamate it with the actual finite configuration,
identifying the copy of $S$.  The domination conclusion of
Proposition~\ref{iso:prime} makes $x$ free over this model from $M_0$.
It is outside the model by Lemma~\ref{geom:nonalgebraic}.  The same
amalgamation can retain any further specified finite entries of the
original configuration.  This is a coupling witnessing that $S$ does
not force $x$.
\end{proof}

\begin{theorem}[The geometry of the minimal type]
\label{geom:pregeometry}
On $P(M)$, for every $M_0\leq_\K M$ of size at least $\lambda_0$,
the preceding finite independence is the independence relation of a
finitary pregeometry.  Its closure, denoted by $\cl_M$, has the
following properties.
\begin{enumerate}
\item For finite $F\subseteq P(M)$ and any maximal independent
      $S\subseteq F$, the set $\cl_M(F)$ consists exactly of those
      points that do not independently extend $S$.  Equivalently,
      $\cl_M(F)$ is the set of points forced by $F$.
\item Strong embeddings over $M_0$ preserve and reflect finite
      closure.  If $M_0\leq_\K A\leq_\K M$, then
      $P(A)$ is relatively closed in $P(M)$.
\item For finite $F$, $|\cl_M(F)|\leq\lambda_0$.  For every infinite
      $F\subseteq P(M)$,
      \[
      |\cl_M(F)|\leq |F|+\lambda_0.
      \]
\end{enumerate}
In particular bases exist, all bases have the same cardinality, and
each point of the span of a basis has finite basis support.
\end{theorem}

\begin{proof}
Choose a maximal independent subset $S$ of a finite $F$.  Every
$y\in F\setminus S$ fails to independently extend $S$, so
Lemma~\ref{geom:containment-test} says that every model containing $S$
in the actual configuration must contain all of $F$.  A point that
fails to independently extend $S$ is therefore forced by $F$.
If $x$ independently extends $S$, the avoiding dominated model from
Lemma~\ref{geom:containment-test} can be adjoined while retaining
$F,x$.  It contains $F$ by the preceding observation and omits $x$.
Thus $F$ does not force $x$.  This proves the first assertion and its
independence of the choice of $S$.

Define closure on arbitrary subsets by taking the union of closures
of finite subsets.  Extensivity and monotonicity hold directly from
the containment test.  To verify idempotence on finite configurations,
suppose that finitely many $y_j$ are in $\cl_M(F)$ and that $x$
independently extends a maximal independent $S\subseteq F$.  The
avoiding model over $S$ can be coupled with the whole finite
configuration $F,(y_j),x$.  It contains $F$ and every $y_j$, because
all these points are forced by $S$, and it omits $x$.  Hence the
$y_j$ cannot force $x$.  Finite character now yields idempotence for
all subsets.

For exchange, it suffices by finite character to take finite $F$ and
points $x,y$ with
\[
y\in\cl_M(F\cup\{x\})\setminus\cl_M(F).
\]
If $x\in\cl_M(F)$, idempotence would imply $y\in\cl_M(F)$, so
$x\notin\cl_M(F)$.  Let $S$ be a maximal independent subset of $F$.
Then $S,x$ and $S,y$ are independent.  Since every member of $F$ is
already forced by $S$, the set $S\cup\{x\}$ is maximal independent
in $F\cup\{x\}$.  The displayed membership says that $S,x,y$ is
dependent.  Permutation invariance of products makes $S,y,x$
dependent as well.  Applying the first assertion with $S\cup\{y\}$
gives $x\in\cl_M(F\cup\{y\})$.  This is exchange.

Preservation and reflection follow from the invariant finite-type
definition and the containment test.  For relative closedness, take
finite $F\subseteq P(A)$ and $x\in\cl_M(F)$.  Choose
$D\leq_\K A$ of size $\lambda_0$ containing $M_0\cup F$, and then
choose a section of $M$ containing $D\cup\{x\}$.  Forced containment
gives $x\in D\subseteq A$.  The finite closure of $F$ lies in this
same $D$, proving the finite cardinal bound.  There are $|F|$ finite
subsets of an infinite $F$, so finite character proves the displayed
bound.

For completeness, a maximal independent subset of a set can be
constructed by well-ordering it and retaining a point precisely when
it is outside the span of the points retained earlier.  Finite
character makes the resulting set independent, and maximality makes
it spanning.  If an independent set is spanned by a finite set of
$n$ points, it has at most $n$ points: successively insert its points
into a spanning list, using exchange to remove one of the original
points whenever a new independent point is inserted.  After $n$
insertions no further independent point is in the span.  This also
proves equality of finite basis sizes.

For infinite bases $H,J$ of the same closed set, choose for each
$j\in J$ a finite support $S_j\subseteq H$.  Then
$\cl(J)\subseteq\cl(\bigcup_{j\in J}S_j)$.  Independence of $H$
implies $H\subseteq\bigcup_{j\in J}S_j$, because a member of $H$
cannot belong to the span of the other members.  Hence
$|H|\leq|J|$; interchanging the bases gives the reverse inequality.
The finite argument excludes the possibility of one basis finite
and the other infinite.  Finally, for a closed $C$, the operator
$U\mapsto\cl(C\cup U)\setminus C$ on $P(M)\setminus C$ satisfies
the same closure and exchange axioms.  Finite character follows by
discarding all but the finitely many entries of $U$ in a dependence
witness.  Applying the preceding constructions to this operator
gives relative bases and their finite supports.
\end{proof}

\subsection{Pointed levels and change of base}

Choose an enumeration $\bar m_0$ of $M_0$ and place it as a strong
model in a sufficiently large well-ordered presentation of $\K$.
Such a placement exists by the order construction and categoricity at
$\lambda_0$.  Fix the stabilized diagram $\Delta$ of this entire
enumeration.  The initial supporting positions of this placement will
be held fixed in later order presentations.

\begin{proposition}[Pointed working levels]
\label{geom:pointed-levels}
At all sufficiently closed levels $\mu$ above the bounds for
$|\bar m_0|=\lambda_0$, models with a distinguished strong copy of
$M_0$ whose enumeration has diagram $\Delta$ form a categorical
working scheme.  Arrows preserve the indicated copy pointwise.
For a pointed model $B$ at such a level, $p_0$ has a unique
nonalgebraic extension $p^B$ over $B$; it is its free lift.  This family
lifts freely along pointed arrows, including the permitted mixed
arrows.  Naming further bounded data of a fixed diagram within the
base does not change type computation over the whole base.
\end{proposition}

\begin{proof}
The enumeration is a bounded tuple at $\mu$.  Bounded-diagram
homogeneity in Proposition~\ref{lev:closure-categoricity} carries it
between any two instances of $\Delta$ by an isomorphism of the entire
models.  The pointed version of the working package follows by
reserving its supporting positions and including the enumeration in
every relevant diagram, as in
Proposition~\ref{lev:named-package}.

Existence of $p^B$ is mixed free extension of $p_0$, and
Lemma~\ref{geom:nonalgebraic} gives nonalgebraicity.  To prove
uniqueness, take two nonalgebraic extensions over $B$.  A distinguishing
test for singleton types may be chosen with a fixed bound below
$\lambda_0$.  For any such parameter set $Z\subseteq B$, choose
$M_0\leq_\K A\leq_\K B$ of size $\lambda_0$ containing $Z$.
Both types restrict to nonalgebraic extensions of $p_0$ over $A$ and
are equal there by Proposition~\ref{geom:minimal-type}.  Their
diagrams on $Z$ are consequently equal.  The bounded testing theorem
gives equality over $B$.  Uniqueness and transitivity of free
extension prove the assertion for pointed arrows.

All distinguished data here already belong to the parameter model.
An amalgam over that model fixes them, and an identification of types
over that model respects them.  Thus adding them as constants does
not change this type computation.
\end{proof}

\begin{proposition}[Change of base for the geometry]
\label{geom:base-change}
Let $B$ be a pointed model at a sufficiently closed level much larger
than $\lambda_0$.  A finite list of points in an extension of $B$ is
independent over the closed set $P(B)$ in the geometry of
Theorem~\ref{geom:pregeometry} if and only if its type over $B$ is the
product of the corresponding copies of $p^B$.
\end{proposition}

\begin{proof}
A product can be witnessed by successive free realizations in models
containing all preceding entries.  Each realization is outside that
model by Lemma~\ref{geom:nonalgebraic}.  Its point set is relatively
closed, so the realization is outside the geometric span of $P(B)$
and the earlier entries.  This proves one implication.

For the converse let $x$ be a finite list independent over $P(B)$.
Fix a parameter test $Z\subseteq B$ large enough to be used in the
uniform comparison of the type of $x$ over $B$, but with
$|Z|<\lambda_0$.  Choose a regular infinite $\theta<\lambda_0$
with $|Z|<\theta$.  We may and do choose $\lambda_0$ closed enough
that the comparisons and local-character bounds used here fit below
it.

Inside $B$ choose lists $x_i$, $i<\theta$, so that every
$D^\K(x_i,Z)$ equals $D^\K(x,Z)$ and all the $x_i$ together are
independent over $M_0$.  To make the successor choice, include in the
bounded match the entire enumeration of $M_0$, $Z$, and all earlier
copies.  The original $x$, being independent over $P(B)$, is independent
from those earlier copies.  Saturation realizes its joint diagram
inside $B$.  Diagram equality is witnessed by an isomorphism, so it
preserves every finite independence comparison over $M_0$.  The total
parameters at every choice are bounded at the level of $B$.

Take a strong section $N\leq_\K B$ of size $\lambda_0$ containing
$M_0$, $Z$, and all the copies.  Their finite sublists have product
type over $M_0$.  The independent-list and finite-domination
conclusions of Theorem~\ref{iso:independent-lists} give a continuous
witness chain $(A_i:i\leq\theta)$ inside $N$, starting at $M_0$, in
which $x_i$ is free over $A_i$ from $M_0$ and is included in
$A_{i+1}$.  One may keep the chain inside $N$: at each step primeness
embeds the required all-dominated increment into $N$ over its already
placed base and the specified next list.  At limits finite domination
keeps the next finite list free, and unions remain strong in $N$.

Since $\cf(\theta)>|Z|$, local character applies to $\tp(Z/A_\theta)$.
It is free from some $A_i$.  For a later $j<\theta$, its restriction
to $A_{j+1}$ is free over $A_j$, so $Z$ and $x_j$ have product over
$A_j$.  The type of $x_j$ over $A_j$ is free from $M_0$.
The base-change identity in Proposition~\ref{calc:products} therefore
says that $x_j$ has product with $Z$ over $M_0$.  In particular
$D^\K(x_j,Z)$ is determined by $p_0$ and $\tp(Z/M_0)$.

A known product of the $p^B$ has this same joint diagram with $Z$:
products of free lifts are free lifts of products, and one may
restrict the resulting type to $M_0$ together with $Z$.  Since
$D^\K(x,Z)=D^\K(x_j,Z)$, the two candidate types over $B$ agree on
every test of the uniform comparison size.  Bounded testing makes
their full types equal.  This proves the converse.
\end{proof}

\subsection{Eliminating proper extensions with unchanged geometry}

The obstruction considered here is related to Vaughtian pairs: proper
extensions of the same cardinality with no new realizations of a fixed
nonalgebraic type; compare \cite{GrossbergVanDieren2006}, where the type
is based on a limit model.  Here finite marked systems and exact boundary
lifting eliminate the obstruction in the pointed setting.

The next use of the marked construction requires a type prescription
that can be chosen from a fixed set before choosing the common level.
We record the uniformity explicitly.

\begin{lemma}[Anchoring a singleton type]
\label{geom:anchored-type}
\label{geom:root-prescription}
There is a fixed cardinal $\beta$, depending on the single-model
template, its singleton comparison bound, and the fixed point, with
the following property.  If $A$ is a sufficiently closed pointed
model and $r$ is a nonalgebraic singleton type over $A$, a list
$c\subseteq A$ of length at most $\beta$ can be named so that:
\begin{enumerate}
\item the joint stabilized diagram of $c,\bar m_0$, together with a
      fixed majority prescription, defines a unique nonalgebraic
      singleton type $r^B$ over every sufficiently closed model $B$
      having that named diagram;
\item the prescription at $A$ is $r$, and the family $r^B$ lifts
      freely along arrows preserving the names, including across
      eligible levels;
\item the prescription occurs at arbitrarily large eligible
      single-model levels.
\end{enumerate}
The possible descriptors obtained this way form a set, with their
length bound independent of $\|A\|$.
\end{lemma}

\begin{proof}
Fix a sufficiently large uniform testing arity $\tau$ for singleton
types, including the bound needed to apply compact placement to those
tests.  The approximation construction for
Theorem~\ref{calc:free-extension} gives a list $c=(c_i:i<\beta')$ of
approximations in $A$ that computes $r$ by majority on all these
tests.  Its length and all ancillary approximation data can be bounded
by one cardinal $\beta$, independently of the level.  Choose the
length larger than the exception bound for the tests; increase
$\beta$ also to include the fixed point data.  The prescription is
the value attained by all but the permitted number of entries in
$D^\K(c_i,z)$, for each test $z$ of arity at most $\tau$.

Freely lift the actual $r$ to arbitrarily large eligible models
containing $A$.  Query matching over the fixed approximation list
shows that these lifts compute exactly the same majority
prescription.  They are nonalgebraic by
Lemma~\ref{geom:nonalgebraic}.  For another sufficiently closed
pointed $B$ with the same named diagram, embed $B$ into such a larger
instance, respecting the names.  This is possible by bounded-diagram
homogeneity and mixed placement.  Restrict the lifted type to that
copy of $B$.  It is nonalgebraic, since its realization is outside the
larger parameter model, and query matching makes its bounded answers
the prescribed ones.

Any two types having these answers are equal by the singleton
comparison bound.  Thus the construction is independent of the
larger instance and the chosen embedding.  A free lift along an arrow
preserving the names has the same prescribed answers: the
approximation evaluation is transported by matching each query over
the named list inside the old base.  Uniqueness therefore identifies
it with the prescribed receiving type.  At $A$ the answers are those
of $r$ by the choice of $c$.

The list lengths have a fixed set bound.  Their stabilized diagrams,
the fixed testing arity, and the majority rule range over a set of
descriptors.  Only descriptors realized at arbitrarily high eligible
single-model levels are needed; the preceding construction proves
that the descriptor extracted from an actual $r$ has this property.
The named diagram itself can be witnessed at one fixed eligible
level large enough for arity $\beta$, by saturation for diagrams from
larger levels.  Proposition~\ref{lev:named-package} therefore applies
with a lower threshold depending on $\beta$ and the original working
scheme, rather than on $\|A\|$.  In particular, extracting the names
from a model of size $\lambda$ does not require moving above
$\lambda$ in order to use the pointed working scheme at that level.
\end{proof}

\begin{theorem}[No proper pair with unchanged points at one level]
\label{geom:no-growth}
There is a sufficiently large cardinal $\lambda>\lambda_0$ such that
there is no proper strong inclusion $A<_\K B$ of pointed models of
size $\lambda$ with $P(A)=P(B)$.
\end{theorem}

\begin{proof}
Take the set of anchored descriptors from
Lemma~\ref{geom:anchored-type}.  Choose $\lambda$ at a common level
for all finite marked constructions for these descriptors, using
Theorem~\ref{close:common-levels}.  Require also all the relevant
single-model and pointed bounds, and $\lambda>\lambda_0+\beta$.
This is a choice uniform in the possible descriptor.  In particular
the choice does not depend on a pair that might later be assumed to
exist.

Suppose toward a contradiction that $A<_\K B$ is such a pair, and
choose $a\in B\setminus A$.  Apply
Lemma~\ref{geom:anchored-type} to $r=\tp(a/A)$ and name its list
$c$.  We work in the pointed scheme also naming this list, with its
fixed joint diagram.  Write $\mathcal C_n$ for its $n$-fold marked
iteration with the root prescription $r$, as in
Theorem~\ref{iter:finite-systems}.

Every one-step all-dominated increment of this scheme has no new
$P$-points.  Indeed, pointed categoricity carries its base, including
the named list $c$, onto $A$.  Transporting the witness $A<_\K B$
back gives an extension with a realization of the required $r$ and
unchanged $P$.  Primeness embeds the all-dominated increment into
that extension over its base and mark.  Since strong maps preserve
and reflect $P$ and fix the base, the increment has exactly the base
points.

It follows by induction that the $P$-points in the top of every
$\mathcal C_n$ object are precisely those in its root.  For the
induction step, its lower and upper layers are $\mathcal C_{n-1}$
objects and hence each has the points of its own root.  Projection
onto the new axis is a valid one-step marked increment by
Theorem~\ref{iter:finite-systems}.  The two roots therefore have the
same points, proving the assertion for the entire object.

We next construct coherent finite systems.  For each $n<\omega$
choose an object $J_n\in\mathcal C_n$.  For every
$s\subseteq\{1,\dots,n\}$ choose an isomorphism
\[
e_{n,s}:J_{|s|}\longrightarrow P_sJ_n
\]
such that these maps commute with all further face projections, using
the increasing identifications of finite index sets.  Here the full
face map is the identity.  To make these choices, induct first on
$n$ and then on the size of a proper subset $s$.  For a proper subset
$u$ of the axes of $J_{|s|}$, let $t$ be the corresponding subset of
$s$.  The boundary map for $e_{n,s}$ on this face is
$e_{n,t}e_{|s|,u}^{-1}$, with the indicated index identifications.
The already chosen maps make these boundary prescriptions compatible.
Theorem~\ref{iter:exact-boundary} extends them to the desired
isomorphism.  At dimension zero it is simply pointed categoricity.
This constructs every $e_{n,s}$ with the asserted commutation.

For the finite subsets $s$ of any ordinal $I$, take a tagged copy of
the top vertex of $J_{|s|}$.  If $s\subseteq t$, use the top component
of $e_{|t|,u}$ followed by the inclusion of that face into the top,
where $u$ records the positions of $s$ in the increasing list of $t$.
The preceding commutation gives a directed system of strong maps.
All its roots are copies of the same actual $J_0$, and no vertex has
more $P$-points than that root.

The marks corresponding to different indices of $I$ are distinct in
this directed system.  It suffices to inspect a two-axis object.  Its
new mark lies in the root of its upper one-axis layer, whereas the
old mark is the nonalgebraic mark of that layer and is outside its
root.  Face isomorphisms preserve the marks, so the distinction
persists under all transition maps.

Take $I$ of cardinality $\mu>\lambda$, where $\mu$ is a sufficiently
closed single-model pointed level.  The ordinary directed colimit in
the AEC exists: a coherent diagram of injective maps is realized as
a diagram of inclusions in its direct limit, and the directed-union
axioms apply.  There are $\mu$ finite subsets of $I$, each contributing
at most $\lambda$ elements, and there are $\mu$ distinct marks.  Thus
the colimit $N$ has size exactly $\mu$.  Every element of $N$ comes
from one finite vertex.  Preservation and reflection of $P$ give
\[
P(N)=P(J_0),\qquad |P(N)|\leq\lambda.
\]
The enumeration of its distinguished $M_0$ retains the fixed
stabilized diagram $\Delta$ under the strong colimit embeddings.

On the other hand, there is a pointed model of size $\mu$ with more
than $\lambda$ points.  Start with any pointed model of that size.
Successively realize a nonalgebraic free extension of $p_0$ outside
the current model, for $\lambda^+$ steps.  The receiving models and
their unions may all be kept of size $\mu$.  The new realizations
are distinct and belong to $P$, so the union has at least
$\lambda^+$ points.  Pointed categoricity at $\mu$ would make this
model isomorphic to $N$ over $M_0$, a contradiction.

Only the stationary prescription $r$ was lifted to higher levels.
The assumed witness with unchanged $P$ was used solely at
$\lambda$.  The large colimit used the full AEC directed-union
axioms, rather than a same-level closure assertion for a marked
scheme.
\end{proof}

\begin{corollary}[No-growth and rank on the entire tail]
\label{geom:tail-no-growth}
Increase the threshold $\lambda$ if necessary to include the bounds
for the fixed point diagram.  In arbitrary cardinalities, if
$A\leq_\K B$ are pointed models with $\|A\|\geq\lambda$ and
$P(A)=P(B)$, then $A=B$.
Moreover, for every pointed model $M$ with $\|M\|>\lambda$,
\[
|P(M)|=\rk(P(M))=\|M\|.
\]
The point sets of strong submodels are relatively closed, and
relative bases and finite basis supports are available as in
Theorem~\ref{geom:pregeometry}; no eligibility condition is imposed
on these arbitrary cardinalities.
\end{corollary}

\begin{proof}
Suppose that $a\in B\setminus A$ and $P(A)=P(B)$.  Construct
increasing strong submodels $A_i\leq_\K A$ and $B_i\leq_\K B$ of
size $\lambda$, for $i<\omega$, all containing the distinguished
copy of $M_0$, so that
\[
A_i\leq_\K B_i,\qquad a\in B_i,\qquad
P(B_i)\subseteq A_{i+1}.
\]
For the successor choice, $P(B_i)$ has size at most $\lambda$ and
is contained in $A$, so the L\"owenheim--Skolem axiom places it,
together with $A_i$, in the next $A_{i+1}$.  Choose $B_{i+1}$ to
contain $A_{i+1}\cup B_i\cup\{a\}$.  Coherence makes all the
required inclusions strong.  The two unions $A_\omega\leq_\K
B_\omega$ have size $\lambda$.  The point diagram restricts
correctly to them, and
\[
P(B_\omega)=\bigcup_{i<\omega}P(B_i)
             \subseteq A_\omega.
\]
Thus their point sets are equal, while $a\in B_\omega\setminus
A_\omega$.  This contradicts Theorem~\ref{geom:no-growth}.

Now let $\|M\|=\kappa>\lambda$.  If $|P(M)|<\kappa$, choose a
strong submodel containing $M_0\cup P(M)$ of size at most
$|P(M)|+\lambda_0$.  Enlarge it inside $M$, if necessary, to a
strong submodel of size at least $\lambda$ and still less than
$\kappa$.  It has the same $P$ as $M$, contrary to the assertion
just proved.  Hence $|P(M)|=\kappa$.  If $H$ is a basis of $P(M)$,
the closure bound gives
\[
\kappa=|P(M)|=|\cl_M(H)|\leq |H|+\lambda_0.
\]
Since $\kappa>\lambda\geq\lambda_0$, it follows that
$|H|=\kappa$.  The final assertions were established in
Theorem~\ref{geom:pregeometry} independently of eligibility.
\end{proof}
